% =====================================================================
\chapter{Primitivas e integral de Riemann}
\label{tema:1}
% =====================================================================

El cálculo integral responde a dos preguntas que, en apariencia, no tienen relación: qué función tiene una derivada dada y cuánto vale el área bajo una curva. Este tema plantea las dos por separado. La primera conduce a las primitivas; la segunda, a la integral de Riemann, definida como el número que queda atrapado entre las sumas inferiores y superiores de rectángulos. El teorema fundamental del \cref{tema:2} mostrará que ambas preguntas son la misma. La exposición sigue los capítulos~13 de Spivak~\cite{spivak} y~1 de Apostol~\cite{apostol}, con los ejemplos y ejercicios de las §5.1 y~5.2 de~\cite{openstax1}, de la §5.5 de Strang~\cite{strang} y de la parte~A de la unidad~3 del curso 18.01SC del MIT~\cite{mit1801}.

\section{Primitivas}

\begin{definicion}[Primitiva]\label{def:primitiva}
Sea $I\subseteq\R$ un intervalo y sea $f\colon I\to\R$. Una función $F\colon I\to\R$ es una \emph{primitiva} de $f$ en $I$ si $F$ es derivable en $I$ y $F'(x)=f(x)$ para todo $x\in I$. El conjunto de todas las primitivas de $f$ se denota $\int f(x)\dif x$ y se llama \emph{integral indefinida} de $f$.
\end{definicion}

\begin{proposicion}\label{prop:primitivas-constante}
Sean $F$ y $G$ dos primitivas de $f$ en un intervalo $I$. Existe una constante $C\in\R$ tal que $G(x)=F(x)+C$ para todo $x\in I$. En consecuencia, si $F$ es una primitiva, $\int f(x)\dif x=\{F+C: C\in\R\}$, lo que se escribe abreviadamente $\int f(x)\dif x=F(x)+C$.
\end{proposicion}

La hipótesis de que $I$ sea un intervalo es imprescindible. En $\R\setminus\{0\}$, la función $H(x)=0$ si $x<0$ y $H(x)=1$ si $x>0$ tiene derivada nula y no es constante; por eso las primitivas de $1/x$ en $\R\setminus\{0\}$ son $\ln\abs{x}+C_1$ para $x<0$ y $\ln\abs{x}+C_2$ para $x>0$, con constantes independientes.

Como la derivada es lineal, también lo es la integral indefinida: si $F$ y $G$ son primitivas de $f$ y $g$, entonces $\alpha F+\beta G$ es primitiva de $\alpha f+\beta g$. Leyendo al revés la tabla de derivadas del módulo~0.5 se obtienen las \emph{primitivas inmediatas} de la \cref{tab:inmediatas}; cada fila se comprueba derivando la columna derecha.

\begin{table}[h]
\centering
\small
\begin{tabular}{@{}lll@{}}
\toprule
$f(x)$ & $\int f(x)\dif x$ & Validez \\
\midrule
$x^{r}$, $r\neq-1$ & $\dfrac{x^{r+1}}{r+1}+C$ & $x>0$ (todo $x$ si $r\in\N_0$) \\[4pt]
$\dfrac1x$ & $\ln\abs{x}+C$ & cada intervalo de $\R\setminus\{0\}$ \\[4pt]
$e^{ax}$, $a\neq0$ & $\dfrac{e^{ax}}{a}+C$ & $\R$ \\[4pt]
$\sin(ax)$, $\cos(ax)$, $a\neq0$ & $-\dfrac{\cos(ax)}{a}+C$, \ $\dfrac{\sin(ax)}{a}+C$ & $\R$ \\[4pt]
$\dfrac{1}{\cos^{2}x}=1+\tan^{2}x$ & $\tan x+C$ & cada intervalo donde $\cos x\neq0$ \\[4pt]
$\dfrac{1}{1+x^{2}}$ & $\arctan x+C$ & $\R$ \\[4pt]
$\dfrac{1}{\sqrt{1-x^{2}}}$ & $\arcsin x+C$ & $(-1,1)$ \\[4pt]
$f(ax+b)$, $a\neq0$ & $\dfrac{1}{a}F(ax+b)+C$, con $F'=f$ & donde esté definida \\
\bottomrule
\end{tabular}
\caption{Primitivas inmediatas. La última fila es la regla de la cadena leída al revés.}
\label{tab:inmediatas}
\end{table}

\addtocounter{definicion}{1}%

\section{Particiones y sumas de Darboux}

La idea de la integral es aproximar el área bajo la gráfica de $f$ por rectángulos y refinar la aproximación. Para que la definición sea rigurosa conviene encerrar el área entre una aproximación por defecto y otra por exceso.

\begin{definicion}[Partición]\label{def:particion}
Sea $[a,b]$ un intervalo con $a<b$. Una \emph{partición} de $[a,b]$ es un conjunto finito $P=\{x_0,x_1,\dots,x_n\}$ con $a=x_0<x_1<\dots<x_n=b$. Se escribe $\Delta x_i=x_i-x_{i-1}$. La \emph{norma} de $P$ es $\norm{P}=\max_i\Delta x_i$. Una partición $Q$ es \emph{más fina} que $P$ si $P\subseteq Q$. La \emph{partición uniforme} de $n$ subintervalos es $x_i=a+i\,\frac{b-a}{n}$.
\end{definicion}

\begin{definicion}[Sumas inferior y superior]\label{def:sumas}
Sea $f\colon[a,b]\to\R$ acotada y sea $P$ una partición. Para cada $i$ se definen
$m_i=\inf\{f(x):x\in[x_{i-1},x_i]\}$ y $M_i=\sup\{f(x):x\in[x_{i-1},x_i]\}$, y
\[
\Lsum(f,P)=\sum_{i=1}^{n}m_i\,\Delta x_i,\qquad \Usum(f,P)=\sum_{i=1}^{n}M_i\,\Delta x_i .
\]
\end{definicion}

Si $f\ge0$, $\Lsum(f,P)$ es el área de una unión de rectángulos contenida en la región bajo la gráfica, y $\Usum(f,P)$ el de una unión que la contiene (\cref{fig:darboux}). Siempre $\Lsum(f,P)\le\Usum(f,P)$, porque $m_i\le M_i$.

\addtocounter{figure}{1}%

\begin{lema}[Refinamiento]\label{lem:refinamiento}
Si $Q$ es más fina que $P$, entonces $\Lsum(f,P)\le\Lsum(f,Q)$ y $\Usum(f,Q)\le\Usum(f,P)$.
\end{lema}

\begin{corolario}\label{cor:inferior-superior}
Para dos particiones cualesquiera $P$ y $Q$ de $[a,b]$, $\Lsum(f,P)\le\Usum(f,Q)$.
\end{corolario}

\section{La integral de Riemann}

Por el \cref{cor:inferior-superior}, el conjunto de las sumas inferiores está acotado superiormente por cualquier suma superior. Existen, por tanto, la \emph{integral inferior} $\underline{I}=\sup_P\Lsum(f,P)$ y la \emph{integral superior} $\overline{I}=\inf_P\Usum(f,P)$, y se cumple $\underline{I}\le\overline{I}$.

\begin{definicion}[Integral de Riemann]\label{def:integral}
Una función acotada $f\colon[a,b]\to\R$ es \emph{integrable} (en el sentido de Riemann) si $\underline{I}=\overline{I}$. Ese valor común es la \emph{integral definida} de $f$ en $[a,b]$ y se denota $\int_a^b f(x)\dif x$. Se conviene además que $\int_a^a f=0$ y $\int_b^a f=-\int_a^b f$.
\end{definicion}

La variable $x$ de $\int_a^b f(x)\dif x$ es muda: $\int_a^b f(x)\dif x=\int_a^b f(t)\dif t$. La integral es un número; la integral indefinida, un conjunto de funciones. El teorema fundamental del \cref{tema:2} relacionará ambas.

\begin{teorema}[Criterio de integrabilidad de Riemann]\label{teo:criterio}
Una función acotada $f\colon[a,b]\to\R$ es integrable si y solo si para cada $\varepsilon>0$ existe una partición $P$ con $\Usum(f,P)-\Lsum(f,P)<\varepsilon$.
\end{teorema}

\addtocounter{definicion}{1}%

\addtocounter{definicion}{1}%

\begin{definicion}[Suma de Riemann]\label{def:suma-riemann}
Sea $P$ una partición de $[a,b]$ y sea $\xi=(\xi_1,\dots,\xi_n)$ con $\xi_i\in[x_{i-1},x_i]$. La \emph{suma de Riemann} es $\Rsum(f,P,\xi)=\sum_{i=1}^{n}f(\xi_i)\,\Delta x_i$. Con $\xi_i=x_{i-1}$, $\xi_i=x_i$ o $\xi_i=\frac{x_{i-1}+x_i}{2}$ se obtienen las sumas por la izquierda, por la derecha y del punto medio.
\end{definicion}

\begin{proposicion}\label{prop:sumas-riemann}
Sea $f$ integrable en $[a,b]$ y sea $(P_n)$ una sucesión de particiones con $\Usum(f,P_n)-\Lsum(f,P_n)\to0$. Para cualquier elección de puntos intermedios $\xi^{(n)}$, $\Rsum(f,P_n,\xi^{(n)})\to\int_a^b f$.
\end{proposicion}

La \cref{prop:sumas-riemann} permite calcular límites de sumas reconociéndolas como sumas de Riemann. Por ejemplo, con la partición uniforme de $[0,1]$ y $\xi_k=k/n$,
\[
\lim_{n\to\infty}\frac1n\sum_{k=1}^{n}\Bigl(\frac{k}{n}\Bigr)^{2}=\int_0^1x^{2}\dif x=\frac13 .
\]

\section{Funciones integrables}

\begin{teorema}[Las funciones monótonas son integrables]\label{teo:monotona}
Toda función monótona $f\colon[a,b]\to\R$ es integrable. Además, con la partición uniforme $P_n$, $\Usum(f,P_n)-\Lsum(f,P_n)=\frac{b-a}{n}\abs{f(b)-f(a)}$.
\end{teorema}

\begin{teorema}[Las funciones continuas son integrables]\label{teo:continua}
Toda función continua $f\colon[a,b]\to\R$ es integrable, y para la partición uniforme $\Usum(f,P_n)-\Lsum(f,P_n)\to0$.
\end{teorema}

Las mismas ideas prueban que una función acotada con un número finito de discontinuidades también es integrable, y que su integral no cambia si se modifica su valor en un número finito de puntos (\cref{eje:1-un-punto}).

\section{Propiedades de la integral}

\begin{teorema}[Propiedades de la integral]\label{teo:propiedades}
Sean $f,g$ integrables en $[a,b]$ y sean $\alpha,\beta\in\R$.
\begin{enumerate}
  \item (\emph{Linealidad}) $\alpha f+\beta g$ es integrable y $\int_a^b(\alpha f+\beta g)=\alpha\int_a^b f+\beta\int_a^b g$.
  \item (\emph{Aditividad}) Si $a<c<b$, $f$ es integrable en $[a,c]$ y en $[c,b]$, y $\int_a^b f=\int_a^c f+\int_c^b f$. Con el convenio de la \cref{def:integral}, la igualdad vale para cualesquiera $a,b,c$ en un intervalo donde $f$ sea integrable.
  \item (\emph{Monotonía}) Si $f\le g$ en $[a,b]$, entonces $\int_a^b f\le\int_a^b g$.
  \item (\emph{Acotación}) Si $m\le f(x)\le M$ en $[a,b]$, entonces $m(b-a)\le\int_a^b f\le M(b-a)$.
  \item (\emph{Valor absoluto}) $\abs{f}$ es integrable y $\bigl|\int_a^b f\bigr|\le\int_a^b\abs{f}$.
\end{enumerate}
\end{teorema}

\addtocounter{definicion}{1}%

\addtocounter{definicion}{1}%

\addtocounter{definicion}{1}%

